% =========================
% puzzles/pXX.tex
% Final - The Balancer's Verdict (Condensed)
% =========================

\PuzzleChapterIntro{The Balancer’s Verdict}{Geometry \textbullet\ Decision Chain \textbullet\ Five-Letter Code}{A convex quadrilateral is cut by rails through its diagonal intersection, producing two corner tiles whose areas must satisfy a square-root inequality. The task is to justify a forced chain of geometric reductions and identify the only consistent choice on each of five tumblers, whose initials reveal the final five-letter code.}
w
\Lore{
The Balancer’s Seal is not a lock so much as a verdict.

A quadrilateral is scratched into slate; then two smaller “corner tiles” are cut out by parallels through the diagonal intersection.
The clerk recites the Balance Law as if it were accounting:

\begin{quote}
\textit{The whole is at least the sum of the roots of its corner debts.}
\end{quote}

Five tumblers hang beneath the slate, each offering four rune-words.
Pick one word per tumbler; take their first letters in order; speak the resulting name.
}

\Problem
\scriptsize
A convex glass floor has four corners \(P,Q,R,S\) in that order. The two glass diagonals meet at a pin \(O = PR\cap QS\).

Through \(O\), draw four rails parallel to the sides to define the corner cuts:
Define points \(X \in PQ\) and \(Y \in QR\) such that \(OX \parallel QR\) and \(OY \parallel PQ\).
Similarly, define \(V \in SP\) and \(W \in RS\) such that \(OV \parallel RS\) and \(OW \parallel SP\).

This partitions the glass into four tiles: \(PXOV\), \(QXOY\), \(RYOW\), \(SVOW\),
where \(QXOY\) and \(SVOW\) are parallelograms by construction.

Let \(H=[PQRS]\), \(A=[PXOV]\), and \(B=[RYOW]\), where \([\cdot]\) denotes area.

\medskip
An engraving on the key claims the \textbf{Balance Law}:
\[
\sqrt{H}\ \ge\ \sqrt{A}+\sqrt{B}.
\]

\medskip
The Balancer Key’s five tumblers are labeled I--V. Each tumbler offers four rune-words; exactly one choice per tumbler can be
\emph{justified} from the geometry in a way that makes the Balance Law unavoidable in every valid configuration.
Choose the correct rune-word on each tumbler to form a 5-letter name.

\medskip

\noindent\textbf{Tumbler I (What reshaping is allowed that preserves rails and rescales every area uniformly?).}
\[
\boxed{\textbf{Affine}}\quad\boxed{\textbf{Twist}}\quad\boxed{\textbf{Crumple}}\quad\boxed{\textbf{Break}}
\]

\noindent\textbf{Tumbler II (What may be imposed after such reshaping to simplify the geometry of the diagonals?).}
\[
\boxed{\textbf{X-orthogonalize}}\quad\boxed{\textbf{Parallelogram}}\quad\boxed{\textbf{Encircle}}\quad\boxed{\textbf{Equalize}}
\]

\noindent\textbf{Tumbler III (What hidden alignment appears for the two corner-tiles?).}
\[
\boxed{\textbf{Induce}}\quad\boxed{\textbf{Round}}\quad\boxed{\textbf{Jagged}}\quad\boxed{\textbf{Split}}
\]

\noindent\textbf{Tumbler IV (Which area recipe is correct for a \emph{rifted} 4-corner pane?).}
\[
\boxed{\textbf{Orthocross}}\quad\boxed{\textbf{Heron-root}}\quad\boxed{\textbf{Beam}\times\textbf{Height}}\quad\boxed{\textbf{Diagonal-song}}
\]
Here the recipes mean:
\begin{itemize}
\item \textbf{Orthocross:} for a quadrilateral whose diagonals are perpendicular, with diagonal lengths \(d_1,d_2\),
\[
[\ ]=\frac{d_1d_2}{2}.
\]
\item \textbf{Heron-root:} the triangle formula \(\sqrt{p(p-a)(p-b)(p-c)}\).
\item \textbf{Beam}\(\times\)\textbf{Height:} base times altitude over \(2\) (triangle-only).
\item \textbf{Diagonal-song:} \([\ ]=d_1d_2\) (a tempting song, but missing the factor \(1/2\)).
\end{itemize}

\noindent\textbf{Tumbler V (Which 1-line identity is the stitching thread?).}
\[
\boxed{\textbf{Merge}}\quad\boxed{\textbf{Average}}\quad\boxed{\textbf{Cancel}}\quad\boxed{\textbf{Swap}}
\]


\medskip
\VaultDirective{What 5-letter name appears?}

\SolutionPage

\textbf{I. Choose \textbf{Affine}.} The construction of \(X,Y,V,W\) uses only parallel lines, and any affine map preserves parallelism.
All areas scale by a common factor \(\lambda\), so \(\sqrt{H},\sqrt{A},\sqrt{B}\) all scale by \(\sqrt\lambda\); the inequality is invariant.

\smallskip

\textbf{II. Choose \textbf{X-orthogonalize}.} Apply an affine map fixing \(O\) that sends the two diagonal directions to perpendicular directions (this is possible since the diagonals are non-parallel); hence we may assume
\[
PR\perp QS.
\]

\smallskip

\textbf{III. Choose \textbf{Induce}.} In \(\triangle PQR\), \(O\in PR\) and \(OX\parallel QR\), so \(\triangle POX\sim\triangle PRQ\), giving \(\frac{PX}{PQ}=\frac{PO}{PR}\).
In \(\triangle PRS\), \(OV\parallel RS\) gives \(\triangle POV\sim\triangle PRS\), so \(\frac{PV}{PS}=\frac{PO}{PR}\).
Thus \(\frac{PX}{PQ}=\frac{PV}{PS}\), hence \(XV\parallel QS\) in \(\triangle PQS\).
Similarly, \(OY\parallel PQ\) and \(OW\parallel SP\) yield \(\frac{RY}{RQ}=\frac{RO}{PR}=\frac{RW}{RS}\), so \(YW\parallel QS\) in \(\triangle RQS\).
Since \(QS\perp PR\) and \(PO,RO\subset PR\), we have \(PO\perp XV\) and \(RO\perp YW\).

\smallskip

\textbf{IV. Choose \textbf{Orthocross}.} With erected diagonals,
\[
H=[PQRS]=\frac{PR\cdot QS}{2}.
\]
Also \(A=[PXOV]=\frac{PO\cdot XV}{2}\) and \(B=[RYOW]=\frac{RO\cdot YW}{2}\).
From \(XV\parallel QS\), \(\triangle POV\sim\triangle PRS\) gives \(XV=\frac{PO}{PR}\,QS\); similarly \(YW=\frac{RO}{PR}\,QS\).
Hence
\[
A=\frac{QS}{2PR}PO^2,\qquad B=\frac{QS}{2PR}RO^2.
\]

\smallskip

\textbf{V. Choose \textbf{Merge}.} Therefore
\[
\sqrt{A}+\sqrt{B}=\sqrt{\frac{QS}{2PR}}\,(PO+RO)=\sqrt{\frac{QS}{2PR}}\cdot PR=\sqrt{\frac{PR\cdot QS}{2}}=\sqrt{H},
\]
so equality holds and the forced rune initials spell \(\boxed{\texttt{AXIOM}}\).

\FinalAnswer{AXIOM}

\NextPage